\section{A global odd-lattice localization theorem}
\label{app:odd-lattice-localization}

The following result applies to arbitrary incoming magnitudes. It is a
necessary condition for competitive parity correlation, not the sharp
atomic inequality.

\begin{theorem}[Constant residual energy near an odd lattice]
\label{thm:odd-lattice-localization}
For every $c>0$ there are constants $0<a_-<a_+<\infty$, $K<\infty$,
and $n_c$ such that the following holds for odd $n\ge n_c$.
If $u\in\R^n\setminus\{0\}$ satisfies
\[
 \frac{\left|\E\chi_{[n]}(X)\tanh\left(\sum_i u_iX_i\right)\right|}
      {\|u\|_2}\ge\frac cn,
\]
then there is $a\in[a_-,a_+]$ with
\[
 \sum_{i=1}^n\operatorname{dist}\bigl(|u_i|,(2\mathbb Z+1)a\bigr)^2
 \le K.
\]
Moreover $\|u\|_2$ lies between two positive constant multiples of
$\sqrt n$. If $r_i$ are nearest positive odd integers, so that
$\bigl||u_i|-ar_i\bigr|=\operatorname{dist}(|u_i|,(2\mathbb Z+1)a)$,
then $\sum_i r_i^2=O_c(n)$. All constants are independent of $n$.
\end{theorem}

\begin{proof}
Input sign changes permit assuming $u_i\ge0$. A zero weight makes the
parity coefficient zero, so under the hypothesis all weights are positive.
Write that coefficient as $C_n(u)$ and put
\[
 r_p(a)=\int_0^\infty\frac{|\sin(at)|^p}{\sinh(\pi t/2)}\,dt,
 \qquad B_p=\int_0^\pi\sin(\theta)^p\,d\theta\qquad(p\ge3).
\]
We first prove the uniform comparison
\begin{equation}\label{eq:absolute-diagonal-comparison}
 r_p(a)\le\frac{B_p}{B_3}r_3(a)\qquad(a>0,\ p\ge3).
\end{equation}
After the substitution $\theta=at$ and folding consecutive periods, the
kernel on $(0,\pi)$ is
\[
 K_a(\theta)=\sum_{j\ge0}
 \frac1{\sinh\bigl(\pi(\theta+j\pi)/(2a)\bigr)}.
\]
It is decreasing and convex. On $(0,\pi/2)$ its symmetrization
$L_a(\theta)=K_a(\theta)+K_a(\pi-\theta)$ is decreasing, because
$L_a'(\theta)=K_a'(\theta)-K_a'(\pi-\theta)\le0$.
Under the probability measure proportional to $\sin^3\theta\,d\theta$
on this interval, $\sin^{p-3}\theta$ is increasing and $L_a$ is
decreasing. Their covariance is therefore nonpositive. Dividing the
resulting integral inequality by the positive integral of $L_a$ gives
\eqref{eq:absolute-diagonal-comparison}. The inverse-linear endpoint
singularity is integrable against $\sin^3\theta$, so truncation justifies
the covariance argument directly.

The beta integral implies $\sup_{p\ge3}\sqrt p\,B_p<\infty$. Hence,
for an absolute constant $A_0$,
\[
 g_n(a):=\frac{\sqrt n\,r_n(a)}a
 \le H(a):=A_0\frac{r_3(a)}a.
\]
The continuous function $H$ is bounded and tends to zero as $a\downarrow0$
and as $a\to\infty$. Indeed $|\sin(at)|^3\le a^3t^3$ gives
$r_3(a)/a=O(a^2)$ at zero. At infinity, splitting the integral at $1/a$
and $1$, using $1/\sinh(\pi t/2)\le2/(\pi t)$ near zero and exponential
decay beyond one, gives $r_3(a)=O(1+\log a)$.

The exact sine representation of tanh, H\"older's inequality, and the
geometric-mean--root-mean-square inequality give
\[
 |C_n(u)|\le\prod_i r_n(u_i)^{1/n},\qquad
 \frac{n|C_n(u)|}{\|u\|_2}\le\prod_i g_n(u_i)^{1/n}.
\]
Choose a bound $G>c$ for $H$ and set $\delta=c^2/(2G)$. If fewer than
$n/2$ of the values $g_n(u_i)$ were at least $\delta$, their geometric
mean would be at most $\sqrt{G\delta}=c/\sqrt2$, a contradiction.
Thus at least $n/2$ weights lie in a fixed compact interval
$[b,B]\subset(0,\infty)$ containing $\{a:H(a)\ge\delta\}$.

A coordinatewise nondecreasing function from the cube to $[-1,1]$ has
absolute top Fourier coefficient at most
\[
 \kappa_n=2^{1-n}\binom{n-1}{(n-1)/2}=O(n^{-1/2}).
\]
For completeness, average over coordinate permutations and express the
resulting increasing sequence on Hamming levels as a convex combination
of level thresholds and a constant. The alternating binomial identity
bounds each threshold coefficient by the displayed central value.
Applying this to the tanh ridge and using its assumed gain gives
$\|u\|_2\le A\sqrt n$ for a constant $A=A(c)$. The moderate-coordinate
count gives $\|u\|_2\ge b\sqrt{n/2}$, and therefore
\begin{equation}\label{eq:lattice-correlation-lower}
 |C_n(u)|\ge\frac\beta{\sqrt n},\qquad \beta=cb/\sqrt2>0.
\end{equation}

Put $P(t)=\prod_i|\sin(u_it)|$ and $k(t)=1/\sinh(\pi t/2)$.
Let $J$ index the moderate coordinates and write $N=|J|\ge n/2$.
On every interval $I$ of length one, uniformly in its location, for
$v\in[b,B]$ and $p\ge1$,
\[
 \int_I|\sin(vt)|^p\,dt\le\frac{C(b,B)}{\sqrt p}.
\]
This follows by substitution and covering the resulting interval by a
bounded number of full or partial periods, using $B_p=O(p^{-1/2})$.
H\"older over the coordinates in $J$ consequently gives
\begin{equation}\label{eq:lattice-unit-interval}
 \int_I P(t)\,dt\le\frac{C_1}{\sqrt n},\qquad
 \int_I P(t)^{1/2}\,dt\le\frac{C_2}{\sqrt n}.
\end{equation}
The second bound uses exponent $N/2$ in the single-coordinate integrals.
For sufficiently large $n$ this exponent is at least one.

Summing unit intervals and using exponential decay of $k$ shows that
\[
 \int_T^\infty P(t)k(t)\,dt
 \le\frac{C_3e^{-\pi T/2}}{\sqrt n}\qquad(T\ge1).
\]
Choose a fixed $T$ making this at most $\beta/(4\sqrt n)$.
Near zero, arithmetic--geometric mean gives
\[
 P(t)\le t^n\prod_i u_i\le(At)^n.
\]
Choose a fixed $0<t_0<\min\{1,1/(2A)\}$. Since $k(t)\le2/(\pi t)$,
the integral on $(0,t_0)$ is exponentially small in $n$, and is at most
$\beta/(4\sqrt n)$ for large $n$. The sine representation and
\eqref{eq:lattice-correlation-lower} now imply
\[
 \int_{t_0}^T P(t)k(t)\,dt\ge\frac\beta{2\sqrt n}.
\]

Let $M=\max_{[t_0,T]}P$. Pointwise $P\le\sqrt M\,P^{1/2}$.
Because $k$ is bounded on this compact interval, the second estimate in
\eqref{eq:lattice-unit-interval} gives
\[
 \frac\beta{2\sqrt n}
 \le\int_{t_0}^T Pk\le\frac{C_4\sqrt M}{\sqrt n}.
\]
Thus $M\ge\mu:=(\beta/(2C_4))^2>0$, independently of $n$.
Choose $t_*\in[t_0,T]$ with $P(t_*)\ge\mu$, and let
$\rho_i=\operatorname{dist}(u_it_*,\pi/2+\pi\mathbb Z)\in[0,\pi/2]$.
Since $-\log\cos\rho\ge\rho^2/2$,
\[
 \sum_i\rho_i^2\le-2\log P(t_*)\le-2\log\mu.
\]
Taking $a=\pi/(2t_*)\in[\pi/(2T),\pi/(2t_0)]$ and dividing by
$t_*^2$ proves the asserted constant residual bound. Finally,
\[
 a^2\sum_i r_i^2\le2\|u\|_2^2+2K=O_c(n),
\]
which gives the remaining conclusion.
\end{proof}

The theorem does not force the residual energy to vanish, nor does it
bound the number of nontrivial odd harmonics. A sharp global atom theorem
requires a uniform signed comparison inside this localized family.
